\documentclass[12pt,a4paper]{article}

% ══════════════════════════════════��═══════════════════════════
% PACKAGES
% ══════════════════════════════════════════════════════════════
\usepackage[utf8]{inputenc}
\usepackage[T1]{fontenc}
\usepackage{lmodern}
\setlength{\headheight}{14pt}
\addtolength{\topmargin}{-2pt}
\usepackage[margin=2.5cm]{geometry}
\usepackage{amsmath,amssymb,amsthm}
\usepackage{mathtools}
\usepackage{graphicx}
\usepackage{xcolor}
\usepackage{booktabs}
\usepackage{array}
\usepackage{multirow}
\usepackage{longtable}
\usepackage{enumitem}
\usepackage{fancyhdr}
\usepackage{titlesec}
\usepackage{parskip}
\usepackage{tcolorbox}
\tcbuselibrary{skins,breakable}
\usepackage{hyperref}

% ══════════════════════════════════════════════════════════════
% COLOURS
% ══════════════════════════════════════════════════════════════
\definecolor{headerblue}{RGB}{20,60,140}
\definecolor{claimbox}{RGB}{20,60,140}
\definecolor{warnbox}{RGB}{160,30,30}
\definecolor{lightgray}{RGB}{245,245,245}
\definecolor{midgray}{RGB}{200,200,200}
\definecolor{darkgray}{RGB}{80,80,80}
\definecolor{forestgreen}{RGB}{0,110,55}
\definecolor{amber}{RGB}{175,95,0}
\definecolor{teal}{RGB}{0,100,100}
\definecolor{plum}{RGB}{100,20,120}

% ══════════════════════════════════════════════════════════════
% SECTION FORMAT
% ══════════════════════════════════════════════════════════════
\titleformat{\section}
  {\large\bfseries\color{headerblue}}
  {\thesection.}{0.6em}{}
  [\vspace{-0.3em}\textcolor{headerblue}{\rule{\linewidth}{0.8pt}}]

\titleformat{\subsection}
  {\normalsize\bfseries\color{headerblue}}
  {\thesubsection.}{0.5em}{}

\titleformat{\subsubsection}
  {\normalsize\bfseries\color{darkgray}}
  {\thesubsubsection.}{0.5em}{}

% ══════════════════════════════════════════════════════════════
% HEADER / FOOTER
% ══════════════════════════════════════════════════════════════
\pagestyle{fancy}
\fancyhf{}
\fancyhead[L]{\small\color{darkgray}Lawson --- Alzheimer's Disease as Attractor Medicine}
\fancyhead[R]{\small\color{darkgray}OrganismCore \texorpdfstring{$\cdot$}{.} 2026-03-09}
\fancyfoot[C]{\small\thepage}
\renewcommand{\headrulewidth}{0.4pt}

% ══════════════════════════════════════════════════════════════
% CUSTOM ENVIRONMENTS
% ═════════════════════════════════════════════════════════���════
\newenvironment{claimenv}{%
  \begin{tcolorbox}[colback=claimbox,colframe=claimbox,coltext=white,
    boxrule=0pt,arc=3pt,left=8pt,right=8pt,top=6pt,bottom=6pt,breakable]
  \bfseries\centering}{%
  \end{tcolorbox}}

\newenvironment{warnenv}{%
  \begin{tcolorbox}[colback=red!5,colframe=warnbox,boxrule=1pt,
    arc=3pt,left=8pt,right=8pt,top=6pt,bottom=6pt,breakable]}{%
  \end{tcolorbox}}

\newenvironment{derivenv}{%
  \begin{tcolorbox}[colback=lightgray,colframe=midgray,boxrule=0.5pt,
    arc=3pt,left=8pt,right=8pt,top=6pt,bottom=6pt,breakable]}{%
  \end{tcolorbox}}

\newenvironment{predenv}[1]{%
  \begin{tcolorbox}[colback=blue!4,colframe=headerblue,boxrule=1pt,
    arc=3pt,title={\bfseries\color{white}#1},colbacktitle=headerblue,
    left=8pt,right=8pt,top=6pt,bottom=6pt,breakable]}{%
  \end{tcolorbox}}

\newenvironment{memenv}[1]{%
  \begin{tcolorbox}[colback=teal!5,colframe=teal,boxrule=1pt,
    arc=3pt,title={\bfseries\color{white}#1},colbacktitle=teal,
    left=8pt,right=8pt,top=6pt,bottom=6pt,breakable]}{%
  \end{tcolorbox}}

\newenvironment{cellenv}[1]{%
  \begin{tcolorbox}[colback=plum!5,colframe=plum,boxrule=1pt,
    arc=3pt,title={\bfseries\color{white}#1},colbacktitle=plum,
    left=8pt,right=8pt,top=6pt,bottom=6pt,breakable]}{%
  \end{tcolorbox}}

% ── Theorem environments ──────────────────────────────────────
\newtheorem{theorem}{Theorem}
\newtheorem{axiom}{Axiom}
\newtheorem{corollary}{Corollary}
\newtheorem{definition}{Definition}
\newtheorem{prediction}{Prediction}

% ══════════════════════════════════════════════════════════════
% HYPERREF
% ══════════════════════════════════════════════════════════════
\hypersetup{
  colorlinks=true, linkcolor=headerblue,
  citecolor=headerblue, urlcolor=headerblue,
  pdftitle={Alzheimer's Disease as Attractor Medicine},
  pdfauthor={Eric Robert Lawson},
  pdfsubject={Identity Axis Framework --- Alzheimer's Application},
  pdfkeywords={Alzheimer's disease, attractor geometry, MEF2C,
    DYRK1A, REST, engram, memory identity, CREB, BDNF, TrkB,
    neurodegeneration, identity drift, OrganismCore,
    tau, neuronal identity, epigenetic}
}

% ══════════════════════════════════════════════════════════════
% MACROS
% ══════════════════════════════════════════════════════════════
\newcommand{\RAD}{R_{\mathrm{AD}}}
\newcommand{\RMEM}{R_{\mathrm{mem}}}
\newcommand{\IA}{\mathcal{A}}
\newcommand{\Hub}{\mathcal{H}}

% ══════════════════════════════════════════════════════════════
\begin{document}
% ══════════════════════════════════════════════════════════════

% ── TITLE BLOCK ───────────────────────────────────────────────
\begin{center}
  {\LARGE\bfseries\color{headerblue}
    Alzheimer's Disease as Attractor Medicine}\\[0.6em]

  {\large
    On the Two Simultaneous Identity Drift Events\\
    in Neurodegeneration, Their Geometric Structure,\\
    Their Drug Targets, and the Difference Between\\[0.3em]
    \textit{Neuronal Identity and Memory Identity}}\\[1.0em]

  {\normalsize Eric Robert Lawson}\\[0.2em]
  {\small OrganismCore (Independent Research)\\
    \href{mailto:OrganismCore@proton.me}{OrganismCore@proton.me}
    \quad$\cdot$\quad
    ORCID: \href{https://orcid.org/0009-0002-0414-6544}
      {0009-0002-0414-6544}}\\[0.5em]

  {\small 2026-03-09}\\[0.7em]
  \rule{0.6\linewidth}{0.4pt}\\[0.5em]

  {\footnotesize
    \textit{Framework DOI:}\;
    \href{https://doi.org/10.5281/zenodo.18898788}
      {10.5281/zenodo.18898788}\\[0.2em]
    \textit{Repository:}\;
    \href{https://github.com/Eric-Robert-Lawson/attractor-oncology}
      {\texttt{github.com/Eric-Robert-Lawson/attractor-oncology}}
  }\\[0.5em]
  \rule{0.6\linewidth}{0.4pt}
\end{center}

\vspace{0.5em}

% ── CENTRAL CLAIM BOX ─────────────────────────────────────────
\begin{claimenv}
Alzheimer's disease is two simultaneous identity drift events,
not one. The neurons are losing their neuronal identity
(MEF2C/NEUROD2 falling, DYRK1A rising, REST sequestered
by tau aggregates, cells attempting aberrant cell cycle
re-entry). The engram cells are simultaneously losing their
memory identity (CREB/BDNF/TrkB falling, synaptic
H3K27ac eroding, the informational structure of
specific memories becoming inaccessible). These have
different mechanisms, different drug targets, and
different intervention windows. The framework resolves both.
The neuron can be rescued. The information may not be
lost --- only inaccessible.
\end{claimenv}

\vspace{1em}

% ══════════════════════════════════════════════════════════════
\begin{abstract}
\noindent
We apply the Identity Axis Framework to Alzheimer's disease,
identifying it as the first framework entry with two
simultaneous and mechanistically distinct identity drift
events operating at two different organisational levels.

\textbf{Drift Event~1 (Cellular):} Cortical and hippocampal
neurons lose their committed neuronal identity. The Identity
Anchor is MEF2C\,/\,NEUROD2; the Convergence Hub is DYRK1A.
DYRK1A phosphorylates and inactivates MEF2C (Identity
Anchor suppression) while simultaneously phosphorylating
tau (accelerating aggregation) and promoting aberrant cell
cycle re-entry. REST, which normally acts as an epigenetic
guardian of the correct neuronal attractor in aging,
is sequestered in autophagosomes by tau aggregates,
removing the checkpoint. The identity axis ratio is
$\RAD = \text{MEF2C} / \text{DYRK1A}$; the drug target
derived from the Hub is a CNS-penetrant DYRK1A inhibitor.
Leucettine~L41 and GNF2133 are in trials.

\textbf{Drift Event~2 (Memory):} Engram cells lose their
memory identity --- the specific synaptic and epigenetic
signature that encodes a particular memory. The Identity
Anchor is the CREB/BDNF/TrkB axis; the Hub is HDAC2/HDAC3
over-activity and histone deacetylation of H3K27ac at
memory-encoding loci. The memory is not erased but
\emph{inaccessible}: the informational structure has drifted
below retrieval threshold. The drug target is HDAC2/3
inhibition; TrkB agonism; and targeted
chromatin-remodelling approaches that restore H3K27ac
at engram loci without global genomic effects.

We derive the geometric sub-type structure (shallow vs.\
deep $\RAD$), the connection between the Down syndrome
trisomy~21/DYRK1A overexpression and accelerated AD onset,
the relationship between cognitive reserve and high $\RAD$,
and the structural distinction between \emph{lost} memory
and \emph{inaccessible} memory with implications for
therapeutic intervention timing. A third operational
level --- the sleep architecture dimension --- is
derived as an addendum (Section~\ref{sec:sleep}):
it identifies a third hub ($\mathcal{H}_{\mathrm{sub}}$
= glymphatic failure), a third positive feedback loop
(thalamic tau degeneration $\to$ spindle loss $\to$
maintenance collapse), revises the four-component
protocol to a five-component protocol, and adds
sleep spindle density as a non-invasive continuous
biomarker of $\RAD$ decline rate. Twelve falsifiable
predictions are stated for the record, timestamped
2026-03-09.

\medskip\noindent\textbf{Keywords:}
Alzheimer's disease; attractor geometry; MEF2C; DYRK1A;
REST; engram; memory identity; CREB; BDNF; TrkB; tau;
epigenetic drift; neurodegeneration; identity drift;
HDAC2; cognitive reserve; Down syndrome;
OrganismCore; Identity Axis Framework
\end{abstract}

\newpage
\tableofcontents
\newpage

% ══════════════════════════════════════════════════════════════
\section{Preamble: The Two Levels of Identity in
Alzheimer's Disease}
% ══════════════════════════════════════════════════════════════

The author's initial observation, stated before the formal
derivation, was:

\begin{derivenv}
\textit{``Alzheimer's and dementia --- I believe that is
absolutely a matter of identity drift of some sort.
It is a person's memory structure fading.''}

\medskip
--- Eric Robert Lawson, 2026-03-09
\end{derivenv}

This is correct. But more precisely than stated.

\textit{Memory structure fading} captures the clinical
phenomenology but conflates two mechanistically distinct
events that produce the same clinical presentation by
different routes:

\begin{itemize}
  \item A \textbf{person's memory fading} because the
    \emph{neuron that holds the memory has died} --- this is
    the cellular identity drift: the neuron lost its
    identity, re-entered the cell cycle, and died. The
    memory encoded in that neuron is gone because the
    substrate is gone.
  \item A \textbf{person's memory becoming inaccessible}
    because the \emph{informational structure encoding
    the memory has drifted} --- the neuron is alive but
    the synaptic and epigenetic signature that constitutes
    the memory has eroded below retrieval threshold.
    The information may be recoverable.
\end{itemize}

These are the same outcome observationally. They are different
mechanisms geometrically. The distinction matters for
treatment: one is irreversible (dead neurons); the other may
be reversible (inaccessible but not destroyed memories).

\begin{warnenv}
\textbf{The most important clinical implication of the
two-drift model:} The intervention window for Drift
Event~2 (memory identity) is \emph{earlier} than the
intervention window for Drift Event~1 (neuronal death).
A patient in early Alzheimer's disease has neurons
that are drifting but not yet dead. Their memories
may be inaccessible but not yet erased. Intervention
at this stage can target both drifts simultaneously.
A patient in late Alzheimer's disease has lost neurons
irreversibly. Drift Event~1 intervention cannot restore
the dead neurons. Drift Event~2 intervention is the
only remaining option, and it is limited by how much
engram cell substrate survives. This is why early
detection and intervention is not merely preferable ---
it is geometrically necessary.
\end{warnenv}

% ══════════════════════════════════════════════════════════════
\section{Framework Foundations}
% ══════════════════════════════════════════════════════════════

\subsection{The Identity Axis Theorem}

The framework rests on five axioms~\cite{IAT2026}:
\begin{itemize}
  \item Every committed cell occupies a stable attractor
    in the Waddington epigenetic landscape.
  \item Disease is a cell displaced into a false attractor,
    maintained by a convergent Hub.
  \item The Identity Anchor~$\IA$ is the element whose
    activity most directly constitutes what the cell is.
  \item The Convergence Hub~$\Hub$ is the element whose
    activity most directly maintains the displaced state.
  \item The ratio $R = \IA/\Hub$ is a continuous measurable
    coordinate; $\Hub$ is the primary drug target.
\end{itemize}

\subsection{The Inverted REST Geometry: A Framework
Novelty}

Before the two drift events are derived, a structural
novelty must be stated explicitly. In the framework's
treatment of small cell lung cancer (SCLC) subtype
transitions~\cite{IAT2026}, REST was identified as the
Convergence Hub that \emph{silences} neuronal identity
in non-neuroendocrine SCLC cells: REST rising = neuronal
identity lost.

In Alzheimer's disease, the geometry is \emph{inverted}:

\begin{table}[ht!]
\centering
\caption{REST in SCLC vs.\ REST in Alzheimer's disease:
  the structural inversion}
\label{tab:rest_inversion}
\begin{tabular}{p{3.0cm}p{4.5cm}p{4.5cm}}
\toprule
\textbf{Context} &
\textbf{REST role} &
\textbf{Geometric interpretation} \\
\midrule
SCLC (non-NE subtype) &
Hub: rises to SILENCE neuronal identity;
pushes cells to SCLC-Y false attractor &
REST is the disease-maintenance mechanism;
inhibiting REST would restore
neuronal identity \\[6pt]
Healthy aging neuron &
Guardian: rises to PROTECT the correct
neuronal attractor by silencing cell-death,
stress, and cell-cycle genes &
REST rising = protective; loss of REST
removes the checkpoint; REST is NOT the Hub \\[6pt]
Alzheimer's neuron &
Checkpoint LOST: tau aggregates sequester
REST in autophagosomes; nuclear REST
depleted; checkpoint removed &
Loss of REST is a CONSEQUENCE of Hub
activity (DYRK1A driving tau
aggregation), not the Hub itself \\
\bottomrule
\end{tabular}
\end{table}

\noindent\textbf{The framework therefore correctly
identifies} DYRK1A --- not REST --- as the true
Convergence Hub in Alzheimer's disease. REST loss is
the mechanism by which DYRK1A drives the neuron to
the false attractor (via tau-mediated REST sequestration),
but REST loss is downstream of DYRK1A overactivation.
The drug target is DYRK1A. Restoring REST by other means
(independent of tau clearance) would address only the
checkpoint, not the Hub, and would therefore be
insufficient for durable correction.

% ══════════════════════════════════════════════════════════════
\section{Drift Event~1: Cellular Neuronal Identity Loss}
% ══════════════════════════════════════════════════════════════

\subsection{The Identity Anchor: MEF2C and NEUROD2}

\begin{cellenv}{Cellular Identity Anchor --- Drift Event~1}
\textbf{MEF2C} (Myocyte Enhancer Factor 2C):\\
Terminal selector transcription factor for neuronal
identity maintenance in the cortex and hippocampus.
Required not only to \emph{establish} neuronal identity
during development but to \emph{maintain} it throughout
adult life. Loss of MEF2C produces: impaired mitochondrial
function in neurons, elevated oxidative stress,
neurodegenerative phenotype in animal models, and shift
of microglia to the hyperinflammatory/senescent state.
MEF2C variants are genome-wide significant AD risk
loci~\cite{MolNeuro2024,NatImmuno2025}. Loss of MEF2C
is upstream of Alzheimer's pathology, not downstream.

\medskip\noindent\textbf{NEUROD2} (Neuronal Differentiation
Factor 2):\\
Pioneer transcription factor of cortical excitatory
neuron identity. Falls in AD brains by single-nucleus
RNA-seq analysis. Falls at the same chromatin loci
where H3K27me3 is anomalously lost. The correlation
with DYRK1A overactivity is direct.
\end{cellenv}

The neuronal Identity Axis ratio for Drift Event~1:
\begin{equation}
  \RAD = \frac{\text{MEF2C expression} \; / \;
               \text{NEUROD2 activity}}
              {\text{DYRK1A expression} \; / \;
               \text{tau phosphorylation load}}
  \label{eq:RAD}
\end{equation}

\subsection{The Convergence Hub: DYRK1A}

DYRK1A (Dual-specificity tyrosine-Phosphorylation
Regulated Kinase~1A) is the primary Convergence Hub
for neuronal identity loss in Alzheimer's disease.
It is the maintenance mechanism of the false attractor.

\begin{derivenv}
\textbf{DYRK1A as Hub: the convergent mechanisms}\\[6pt]
DYRK1A phosphorylates \textbf{MEF2C} at Ser408 $\to$
MEF2C is inactivated $\to$ Identity Anchor suppressed
$\to$ neuronal identity programme begins to drift.\\[4pt]

DYRK1A phosphorylates \textbf{tau} at multiple sites
(Thr181, Ser202, Thr231, Ser396, Ser404) $\to$ tau
becomes hyperphosphorylated and aggregation-prone $\to$
tau aggregates form neurofibrillary tangles $\to$
tau aggregates sequester REST in autophagosomes $\to$
nuclear REST depleted $\to$ checkpoint removed.\\[4pt]

DYRK1A promotes \textbf{aberrant cell cycle re-entry}
via phosphorylation of LIN52 (MuvB complex) and
CDK4/Rb pathway activation $\to$ post-mitotic neuron
attempts to divide $\to$ abortive mitosis $\to$
apoptosis.\\[4pt]

DYRK1A is overexpressed on chromosome 21 $\to$ in
trisomy~21 (Down syndrome), constitutive DYRK1A
overexpression $\to$ $\RAD$ depressed from birth $\to$
AD onset at median age~40 in Down syndrome individuals.
\emph{This is the geometric derivation of the Down
syndrome/Alzheimer's connection, derived from the Hub
identity of DYRK1A.}
\end{derivenv}

\subsection{The Secondary Hub: Pathological Tau
as Propagation Mechanism}

The tau aggregate serves a dual role in the geometry:
it is both a \emph{product} of DYRK1A Hub activity
(DYRK1A phosphorylates tau, driving aggregation) and
a \emph{spreading mechanism} for the Hub's effect
(tau aggregates spread trans-synaptically, confirmed
by prion-like propagation studies~\cite{DeCalignon2012}).

\begin{derivenv}
\textbf{The propagation cascade:}
\[
\underbrace{\text{DYRK1A} \uparrow}_{\text{Hub activated}}
\;\longrightarrow\;
\underbrace{\text{tau-p} \uparrow}_{\text{phosphorylation}}
\;\longrightarrow\;
\underbrace{\text{tau aggregates}}_{\text{NFT formation}}
\;\longrightarrow\;
\underbrace{\text{REST sequestered}}_{\text{checkpoint lost}}
\;\longrightarrow\;
\]
\[
\longrightarrow\;
\underbrace{\text{cell cycle entry}}_{\text{CDK4/Rb}}
\;\longrightarrow\;
\underbrace{\text{apoptosis}}_{\text{neuronal death}}
\;\longrightarrow\;
\underbrace{\text{tau released}}_{\text{trans-synaptic}}
\;\longrightarrow\;
\underbrace{\text{adjacent neuron}}_{\text{REST sequestered}}
\longrightarrow \cdots
\]
\end{derivenv}

Each neuron receiving tau aggregates loses its nuclear
REST, becomes vulnerable to DYRK1A-driven cell cycle
re-entry, and dies --- releasing more tau to the next
neuron. The Braak staging of AD corresponds directly
to this geometric propagation: the wave front of the
false attractor spreading through the hippocampal
circuit, the entorhinal cortex, and the neocortex~\cite{Braak1991}.

\subsection{The Sub-Type Structure and Cognitive Reserve}

The $\RAD$ framework predicts a continuous subtype
structure, not a binary healthy/diseased classification:

\begin{table}[ht!]
\centering
\caption{$\RAD$ sub-type structure and clinical predictions}
\label{tab:subtypes}
\begin{tabular}{p{2.2cm}p{3.5cm}p{3.0cm}p{3.0cm}}
\toprule
\textbf{$\RAD$ level} &
\textbf{Molecular state} &
\textbf{Clinical phenotype} &
\textbf{Framework prediction} \\
\midrule
High (healthy) &
MEF2C high, DYRK1A low, REST nuclear,
H3K27me3 preserved, tau unmethylated &
Cognitively normal, no clinical symptoms even
if amyloid-PET positive &
Correct attractor. AD pathology present but
not progressing. ``Cognitive reserve.'' \\[6pt]
Intermediate &
MEF2C decreasing, DYRK1A rising,
REST partially nuclear, early tau-p &
MCI (Mild Cognitive Impairment);
early memory complaints &
Near-threshold. Both drifts beginning.
Optimal intervention window. \\[6pt]
Low &
MEF2C low, DYRK1A dominant, REST
depleted, tau aggregates spreading &
Established AD. Progressive cognitive
decline. Hippocampal volume loss. &
False attractor. Neuronal death occurring.
Drift Event~2 also active. \\[6pt]
Floor &
MEF2C absent, DYRK1A fully active,
most hippocampal neurons dead &
Late-stage AD. Loss of ADLs, language,
personality. &
Irreversible substrate loss. Drift Event~1
uncorrectable. Drift Event~2 intervention
limited by surviving engram cell count. \\
\bottomrule
\end{tabular}
\end{table}

\begin{corollary}[Cognitive Reserve as High $\RAD$]
The phenomenon of ``cognitive reserve'' --- individuals
with confirmed amyloid-PET pathology who remain cognitively
intact for years longer than expected --- is predicted
to be the population with measurably higher baseline
$\RAD$. High MEF2C expression in hippocampal and
entorhinal neurons sustains $\RAD$ above threshold
despite amyloid accumulation, preventing the DYRK1A
overactivation cascade. The cognitive reserve
population will show measurably higher MEF2C, lower
DYRK1A phosphorylation of tau, and higher nuclear
REST in accessible biomarker tissue compared to
age-matched individuals with the same amyloid burden
and earlier clinical symptoms.
\end{corollary}

\subsection{The Drug Target: DYRK1A Inhibition}

The Hub identification ($\Hub = \text{DYRK1A}$) directly
yields the primary drug target:

\begin{predenv}{Drug Target: CNS-Penetrant DYRK1A
Inhibitor}
\textbf{Target:} DYRK1A kinase activity inhibition in
cortical and hippocampal neurons.\\[4pt]
\textbf{Mechanism:} DYRK1A inhibition $\to$ MEF2C
phosphorylation reduced $\to$ MEF2C reactivated $\to$
neuronal identity programme restored. Simultaneously:
tau phosphorylation at DYRK1A sites reduced $\to$ tau
aggregation slowed $\to$ REST sequestration slowed
$\to$ REST checkpoint partially restored.\\[4pt]
\textbf{Clinical candidates:}
\begin{itemize}
  \item \textbf{Leucettine L41}: DYRK1A inhibitor,
    Phase~2 clinical trial in Down syndrome/AD
    prevention~\cite{Leucettine2022}. Confirmed CNS
    penetrance. Reduces tau phosphorylation and
    restores synaptic plasticity in animal models.
  \item \textbf{GNF2133}: DYRK1A inhibitor, preclinical
    AD models~\cite{GNF21332023}. Shows tau reduction
    and cognitive improvement (Morris water maze) in
    3xTg-AD mice.
  \item \textbf{EHT~1610}: Selective DYRK1A/CLK1
    inhibitor, neuroprotective in AD cellular models.
\end{itemize}
\textbf{Framework convergence:} The framework derives
DYRK1A as Hub from geometry. The field arrived at
DYRK1A empirically (Down syndrome/AD connection,
tau phosphorylation mapping). They converge on the
same target. This is the same pattern as every
other table entry.
\end{predenv}

% ══════════════════════════════════════════════════════════════
\section{Drift Event~2: Memory Identity Loss}
% ══════════════════════════════════════════════════════════════

\subsection{Memory as an Identity in Waddington Space}

Drift Event~2 is the framework's most novel contribution to
the AD derivation. It extends the concept of cellular identity
to a second level: the identity of what a cell \emph{holds}.

\begin{memenv}{What Memory Identity Means Geometrically}
A memory is encoded as a specific pattern of:
(1)~long-term synaptic potentiation (LTP) at particular
dendritic spines of the engram cell, (2)~epigenetic marks
(H3K27ac, H3K4me3) at the gene loci whose products
maintain those synaptic weights (Arc, Fos, Homer1a,
GluR1/2), (3)~CREB-mediated transcriptional programme
that sustains the synaptic signature over time.

\medskip
Together, these constitute an \emph{attractor state}
for the engram cell at the \emph{informational level}:
a stable configuration that encodes a specific memory.
The memory is accessible so long as the attractor is
stable. The memory becomes inaccessible --- but not
erased --- when the attractor drifts below retrieval
threshold. The identity of \emph{what the cell holds}
has drifted, not the identity of \emph{what the cell is}.
\end{memenv}

\subsection{The Memory Identity Anchor: CREB/BDNF/TrkB}

\begin{equation}
  \IA_{\mathrm{mem}} = \text{CREB}^{\text{(phospho-Ser133)}}
    \;\cdot\; \text{BDNF} \;\cdot\; \text{TrkB}
  \label{eq:anchor_mem}
\end{equation}

\begin{itemize}
  \item \textbf{CREB (Ser133):} Phospho-CREB is the
    master activator of memory-gene transcription.
    CREB phosphorylation at Ser133 drives expression
    of BDNF, Arc, Fos, Homer1a, and other LTP
    maintenance genes. Loss of CREB-Ser133 activity
    means the synaptic weights at engram cell dendritic
    spines are not being actively maintained. The
    synaptic signature erodes. The memory becomes
    inaccessible~\cite{Bhatt2009,Yiu2014}.
  \item \textbf{BDNF:} Brain-Derived Neurotrophic Factor.
    Acts via TrkB to sustain dendritic spine density,
    LTP maintenance, and CREB transcriptional activity
    in hippocampal neurons. BDNF falls early in AD
    (confirmed: entorhinal cortex BDNF reduction
    precedes clinical symptoms)~\cite{Holsinger2000}.
  \item \textbf{TrkB:} BDNF receptor. TrkB signalling
    activates CaMKIV, which phosphorylates CREB,
    completing the memory maintenance loop. Loss of
    TrkB density is directly associated with synaptic
    loss in AD~\cite{Allen2015}.
\end{itemize}

\subsection{The Memory Identity Hub: HDAC2/HDAC3}

\begin{equation}
  \Hub_{\mathrm{mem}} = \text{HDAC2} \;\cdot\; \text{HDAC3}
  \label{eq:hub_mem}
\end{equation}

HDAC2 and HDAC3 deacetylate H3K27ac at memory-encoding
loci (Arc, Fos, GluR1/2, Homer1a, BDNF). Deacetylation
closes the chromatin at these loci $\to$ gene expression
falls $\to$ synaptic maintenance proteins are no longer
produced $\to$ LTP decays $\to$ the memory's epigenetic
signature erodes toward baseline (inaccessible state).

\begin{itemize}
  \item \textbf{HDAC2} is the primary identity axis
    component: elevated in AD hippocampus~\cite{Ganai2016};
    overexpression produces memory impairment in mice;
    specific knockout of HDAC2 restores memory access
    in AD mouse models~\cite{Guan2009}. HDAC2 is the
    molecular lock on the memory attractor.
  \item \textbf{HDAC3} co-participates in the
    deacetylation of the same loci, particularly in
    the consolidation phase of new memory formation.
    HDAC3 inhibition enhances long-term memory
    formation in rodents~\cite{McQuown2011}.
\end{itemize}

The memory identity axis ratio:
\begin{equation}
  \RMEM = \frac{\text{CREB-Ser133 activity} \;\cdot\;
                \text{BDNF}\;\cdot\;\text{TrkB}}
               {\text{HDAC2} \;\cdot\; \text{HDAC3}}
  \label{eq:RMEM}
\end{equation}

When $\RMEM$ falls below retrieval threshold: the memory
becomes inaccessible. The information is preserved in
the dormant synaptic architecture but cannot be retrieved
because the epigenetic signature has drifted below
expression threshold. This is \emph{not} erasure.
It is \emph{inaccessibility}.

\subsection{The Information Preservation Hypothesis}

\begin{warnenv}
\textbf{The most important structural claim of the
two-drift model:} In early and middle Alzheimer's
disease, before the engram neurons themselves die
(Drift Event~1 reaches the engram cell), the memories
are inaccessible, not erased. The synaptic architecture
that encodes the memory is degraded but not fully
destroyed. The epigenetic marks at memory-encoding loci
have drifted below retrieval threshold but the loci
still exist. Restoring $\RMEM$ above retrieval threshold
--- by HDAC2/3 inhibition, TrkB agonism, or CREB
phosphorylation enhancement --- may restore access
to memories that appeared clinically ``lost.''\\[4pt]
\textbf{Experimental confirmation (partial):}
Optogenetic activation of engram cells in AD mouse
models (APP/PS1) that can no longer retrieve memories
behaviorally \emph{restores} memory recall during the
optogenetic stimulation~\cite{Roy2016}. The information
was there. It was inaccessible. The optogenetic
stimulation bypassed the failed retrieval mechanism.
This is the in vivo proof of the Information Preservation
Hypothesis.
\end{warnenv}

\subsection{The Memory Identity Drug Targets}

\begin{predenv}{Drug Targets for Drift Event~2:
Memory Identity Restoration}
\textbf{Target~1: HDAC2 selective inhibitor.}\\
Inhibit HDAC2 specifically at memory-encoding chromatin
loci $\to$ H3K27ac preserved $\to$ Arc, Fos, BDNF,
GluR1/2 expression maintained $\to$ synaptic
maintenance proteins produced $\to$ LTP sustained $\to$
$\RMEM$ raised.\\[6pt]
\textbf{Candidate:} Selective HDAC2 inhibitor (not
pan-HDAC --- pan-HDAC inhibitors such as vorinostat
have cardiovascular toxicity profiles that preclude
chronic use in AD population). Current state: selective
HDAC2 inhibitors in medicinal chemistry development.
No clinical candidate yet approved for AD specifically.
\\[6pt]
\textbf{Target~2: TrkB agonist.}\\
Activate TrkB directly (bypassing BDNF, which does
not adequately cross the blood-brain barrier at
therapeutic doses) $\to$ CaMKIV/CREB phosphorylation
loop restored $\to$ BDNF gene expression re-initiated
$\to$ $\RMEM$ raised.\\[6pt]
\textbf{Candidate:} ANA~12 (TrkB antagonist --- note:
this is the WRONG direction; what is needed is TrkB
agonist). 7,8-DHF (7,8-Dihydroxyflavone): confirmed
TrkB partial agonist, crosses the BBB, improves
cognition in 3xTg-AD mice~\cite{Devi2012}.
LM22A-4: TrkB agonist, preclinical AD
neuroprotection~\cite{Massa2010}.\\[6pt]
\textbf{Target~3: CREB activator.}\\
Directly enhance CREB-Ser133 phosphorylation via
PDE4 inhibition (raises cAMP $\to$ PKA $\to$
CREB-Ser133) or direct PKA activator $\to$ CREB
transcriptional programme restored $\to$ memory
maintenance gene expression reinstated.\\[6pt]
\textbf{Candidate:} Rolipram (PDE4 inhibitor):
improves memory in multiple AD models; tolerability
issues (emesis) in humans but newer selective
PDE4B inhibitors are in development.
\end{predenv}

% ══════════════════════════════════════════════════════════════
\section{The Down Syndrome / Alzheimer's Connection:
A Geometric Derivation}
% ═══════════��══════════════════════════════════════════════════

The connection between trisomy~21 (Down syndrome) and
early-onset Alzheimer's disease is the single clearest
external validation of the DYRK1A Hub identification.

\begin{derivenv}
\textbf{The geometric derivation:}

DYRK1A maps to chromosome 21 (21q22.13). Trisomy~21
produces three copies of DYRK1A instead of two.
Constitutive 1.5× overexpression of DYRK1A from birth.

From the framework:
\[
\RAD = \frac{\text{MEF2C}}{\text{DYRK1A}}
\]

If DYRK1A is constitutively elevated from birth at 1.5×
normal, $\RAD$ is constitutively depressed by
$\approx$33\% below the normal baseline for the entire
lifetime. The neuron spends its entire life closer to
the threshold. Every additional perturbation (aging,
inflammation, stress) that would normally be tolerated
without crossing the threshold crosses it decades
earlier in a DYRK1A-overexpressing neuron.

\medskip
\textbf{Clinical prediction confirmed:}
Virtually all individuals with Down syndrome develop
Alzheimer's-type neuropathology by age~40 and
clinical dementia by age 50~\cite{Hartley2015}.
This is the constitutive $\RAD$ depression predicted
by the geometry. The Hub is overexpressed from
conception. The false attractor is therefore reached
decades earlier.

\medskip
\textbf{Treatment implication:}
In Down syndrome, the Hub overexpression is constitutional
and cannot be reversed by a kinase inhibitor indefinitely.
The correct intervention is lifelong DYRK1A inhibitor
prophylaxis beginning in early adulthood (before the
threshold crossing), not AD treatment after onset.
Leucettine~L41 Phase~2 trial is running in this
population. This is the correct indication derived
from the geometry.
\end{derivenv}

% ══════════════════════════════════════════════════════════════
\section{The Integrated Treatment Protocol}
% ══════════════════════════════════════════════════════════════

\subsection{Architecture: Two Drift Events, Two
Target Sets, One Sequence}

The integrated protocol addresses both drift events
with distinct but temporally overlapping interventions.
The sequence is determined by the geometric dependencies.

\begin{table}[ht!]
\centering
\caption{Integrated treatment protocol by disease stage}
\label{tab:protocol}
\begin{tabular}{p{2.2cm}p{3.8cm}p{3.8cm}p{3.0cm}}
\toprule
\textbf{Stage} &
\textbf{Drift Event~1 intervention} &
\textbf{Drift Event~2 intervention} &
\textbf{Primary goal} \\
\midrule
Pre-symptomatic
(high amyloid-PET,
normal cognition) &
DYRK1A inhibitor (Leucettine L41)
prophylactic &
HDAC2i + TrkB agonist prophylactic &
Prevent threshold crossing on both
drift axes simultaneously \\[6pt]
MCI (early) &
DYRK1A inhibitor (full therapeutic
dose) + MEF2C agonism support &
HDAC2i + TrkB agonist +
PDE4B inhibitor (CREB activation) &
Halt neuronal death; restore memory
identity for accessible memories;
rescue recently inaccessible memories \\[6pt]
Established AD
(moderate) &
DYRK1A inhibitor (slows further loss);
TrkB agonist (neuroprotection) &
HDAC2i (restore access to surviving
engram cells); intensive cholinergic
support &
Slow Drift Event~1 progression;
maximise recovery of inaccessible
memories in surviving neurons \\[6pt]
Late-stage AD &
Neuroprotective support; tau
clearance (lecanemab adjunct) &
HDAC2i where surviving engram
substrate remains &
Quality of life; slow further
irreversible loss \\
\bottomrule
\end{tabular}
\end{table}

\subsection{The Critical Sequence Rule}

\begin{warnenv}
\textbf{The sequence is not arbitrary.}\\[4pt]
Drift Event~2 intervention (HDAC2 inhibition) MUST
be applied while the engram neurons are still alive.
HDAC2 inhibition at dead neurons recovers nothing.
This is why the pre-symptomatic and MCI intervention
windows are geometrically critical: they are the
windows in which the engram cell substrate that
holds the inaccessible memories still exists. Every
month of delayed Drift Event~1 treatment (DYRK1A
inhibition) is months of engram cells dying ---
shrinking the substrate pool available for Drift
Event~2 recovery.
\end{warnenv}

\subsection{Amyloid Therapy Repositioned}

The framework repositions amyloid-targeting therapies
(lecanemab, donanemab) within the geometry:

Amyloid-$\beta$ accumulation is \emph{upstream} of the
Hub cascade (A$\beta$ oligomers activate multiple kinases
including DYRK1A and GSK3$\beta$, which together
hyperphosphorylate tau). Amyloid clearance is therefore
a partial upstream intervention: it reduces one of the
inputs to the Hub cascade without targeting the Hub itself.

\begin{derivenv}
\textbf{Geometric placement of amyloid therapy:}\\
A$\beta$ oligomers $\to$ DYRK1A/GSK3$\beta$ activation
$\to$ tau phosphorylation $\to$ Hub cascade.\\[4pt]
Lecanemab clears A$\beta$ plaques $\to$ reduces
oligomeric A$\beta$ $\to$ reduces one activating
input to the Hub $\to$ partial $\RAD$ improvement.\\[4pt]
This is why lecanemab shows modest but real clinical
effect (29\% slowing of decline in CDR-SB at 18
months~\cite{Lecanemab2023}) but not halting or
reversal: it reduces Hub activation pressure but
does not address the Hub directly. DYRK1A inhibitor
+ lecanemab should produce additive effects because
they target different nodes in the cascade.
\end{derivenv}

% ══════════════════════════════════════════════════════════════
\section{Twelve Falsifiable Predictions for the Record}
% ══════════════════════════════════════════════════════════════

All predictions are derived from the framework
and timestamped 2026-03-09.

\begin{prediction}[Cognitive reserve = high $\RAD$]
\label{pred:reserve}
Individuals with confirmed amyloid-PET positivity and
no cognitive symptoms (``resilient'' population) will
show measurably higher MEF2C expression in accessible
biomarker tissue (CSF or iPSC-derived hippocampal
neurons), lower DYRK1A-mediated tau phosphorylation
(p-tau~181, p-tau~231 in CSF), and higher nuclear REST
levels compared to age-matched individuals with the
same amyloid burden and established MCI or AD symptoms.
The cognitive reserve effect is a high-$\RAD$ effect.
\end{prediction}

\begin{prediction}[DYRK1A inhibitor + lecanemab
additivity]
\label{pred:combo}
In a clinical trial combining a CNS-penetrant DYRK1A
inhibitor (Leucettine L41 or equivalent) with lecanemab
in early AD, the combination arm will show statistically
significantly greater slowing of CDR-SB decline than
either monotherapy arm, because the combination addresses
both the upstream amyloid input (lecanemab) and the
Hub itself (DYRK1A inhibitor). The additive effect
will be largest in the MCI population where both
drugs can reach enough surviving neurons to act on.
\end{prediction}

\begin{prediction}[HDAC2 inhibition rescues
optogenetically inaccessible memories]
\label{pred:hdac2}
In an AD mouse model (APP/PS1 or 3xTg-AD) at the
stage where specific memories are not behaviourally
accessible, HDAC2-selective inhibition will restore
memory retrieval on the fear-conditioning or
object-recognition task at a rate statistically greater
than vehicle control, without requiring optogenetic
stimulation of the engram cell. This is the
drug-mediated equivalent of the Roy~2016 optogenetic
result~\cite{Roy2016}, achieved pharmacologically
by raising $\RMEM$ above retrieval threshold.
\end{prediction}

\begin{prediction}[7,8-DHF + DYRK1A inhibitor
synergy in 3xTg-AD]
\label{pred:trkb}
In 3xTg-AD mice, combined treatment with 7,8-DHF
(TrkB agonist, Drift Event~2 intervention) and
GNF2133 or EHT~1610 (DYRK1A inhibitor, Drift
Event~1 intervention) will show greater improvement
in both synaptic density (Drift Event~2 biomarker)
and tau phosphorylation (Drift Event~1 biomarker)
than either drug alone, consistent with the two-drift
model predicting independent but complementary
mechanisms.
\end{prediction}

\begin{prediction}[$\RAD$ CSF biomarker predicts
conversion rate]
\label{pred:iad}
A composite CSF biomarker panel of MEF2C/DYRK1A ratio
(or equivalent proxy: p-tau~231 as DYRK1A activity
surrogate; NEUROD2 as MEF2C activity surrogate) will
predict conversion from cognitively normal amyloid-positive
to MCI within 3 years with greater accuracy than
amyloid or tau burden alone.
\end{prediction}

\begin{prediction}[Down syndrome prophylactic window
for Leucettine]
\label{pred:ds}
Leucettine~L41 initiated in Down syndrome individuals
between ages 25 and 35 (before the median threshold
crossing) will significantly reduce CSF tau
phosphorylation markers and delay Alzheimer's
symptoms onset by $\geq$5 years compared to untreated
Down syndrome controls, because it directly addresses
the constitutional DYRK1A overexpression that is
the geometric cause of early $\RAD$ threshold
crossing in this population.
\end{prediction}

\begin{prediction}[Selective HDAC2 inhibitor
outperforms pan-HDAC in AD memory models]
\label{pred:selective}
A selective HDAC2 inhibitor with confirmed specificity
over HDAC1 and HDAC3 will show equivalent memory
restoration in an AD behavioural model compared
to non-selective HDAC inhibitors, with significantly
reduced haematological and cardiovascular adverse
effects, confirming the geometric prediction that
HDAC2 (not HDAC1 or the class I/II panel) is the
specific Hub at memory-encoding chromatin loci.
\end{prediction}

\begin{prediction}[REST restoration insufficient
as monotherapy]
\label{pred:rest}
Restoration of nuclear REST levels (by inhibiting
the tau-mediated autophagosomal sequestration of REST,
for example via TFEB activation increasing
autophagosomal clearance) without simultaneous DYRK1A
inhibition will not prevent progression of tau
phosphorylation or neuronal identity loss in an AD
model, because REST restoration addresses the lost
checkpoint but not the active Hub (DYRK1A) that
is continuing to phosphorylate MEF2C and tau
independently of REST status.
\end{prediction}

% ══════════════════════════════════════════════════════════════
\section{The Sleep Architecture Dimension:
The Third Level of the Alzheimer's Derivation}
\label{sec:sleep}
% ══════════════════════════════════════════════════════════════

\begin{warnenv}
\textbf{Addendum --- derived 2026-03-09.}
The original derivation (Sections 1--7) was written
from a single entry direction: the epigenetic and
molecular machinery of cellular and memory identity
drift. It correctly identified two hubs
($\mathcal{H}_{\mathrm{AD}} = \text{DYRK1A}$;
$\mathcal{H}_{\mathrm{mem}} = \text{HDAC2/3}$) and
two identity drift events. It did not identify a
third operational level of the disease --- one that
operates not at the molecular level of individual
cells but at the level of the maintenance architecture
that services both hubs' primary substrates every
night. That third level is sleep. The cross-framework
integration with the Sleep Disorders Attractor
Medicine derivation (same date, same framework,
different entry point) generates three new structural
claims, a revised protocol, and four additional
falsifiable predictions. None were present in the
original derivation. They are added here with full
epistemic precision.
\end{warnenv}

\subsection{The Third Hub: $\mathcal{H}_{\mathrm{sub}}$
= Glymphatic Failure}

The two hubs identified in Sections 3--4 operate at
the molecular level inside or at the surface of
individual neurons. Both require a substrate to act on:
DYRK1A requires extracellular tau fragments to
phosphorylate and propagate; HDAC2/3 require that
memory-encoding chromatin loci remain accessible.
Both substrates are cleared or maintained by the
glymphatic system during NREM slow-wave sleep.

During correctly running NREM sleep, cerebrospinal
fluid pulses through the perivascular space via
slow-oscillation-coupled arterial pulsatility. The
interstitial space expands by $\approx$60\% during
NREM~\cite{Xie2013}, enabling bulk clearance of
beta-amyloid and phosphorylated tau fragments at
approximately twice the rate of wakefulness.

\begin{derivenv}
\textbf{$\mathcal{H}_{\mathrm{sub}}$ as a third hub:}\\[4pt]
Glymphatic failure
$\to$ oA$\beta$ accumulates faster (clearance reduced)
$\to$ $\mathcal{H}_{\mathrm{mem}}$ hub concentration
rises $\to$ $\RMEM$ falls faster.\\[4pt]
Glymphatic failure
$\to$ extracellular tau fragments accumulate faster
$\to$ trans-synaptic tau propagation accelerates
$\to$ $\mathcal{H}_{\mathrm{AD}}$ cascade spreads faster
$\to$ $\RAD$ falls faster.\\[4pt]
\textbf{$\mathcal{H}_{\mathrm{sub}}$ is a hub operating
at the substrate maintenance level}: its failure
accelerates both molecular hubs simultaneously.
\end{derivenv}

This is the geometric reason obstructive sleep apnea
(OSA) is one of the strongest modifiable risk factors
for Alzheimer's disease. OSA interrupts the NREM
maintenance cycle before glymphatic clearance
completes. Each apnoeic arousal restarts the cycle;
the clearance never finishes. CPAP treatment removes
the interruption; the maintenance cycle completes;
$\mathcal{H}_{\mathrm{sub}}$ is partially addressed.

\begin{warnenv}
\textbf{CPAP is not comorbidity management in this
framework.} It is a primary AD treatment component.
It addresses the third hub --- substrate maintenance
failure --- that the original four-component protocol
does not name.
\end{warnenv}

\subsection{The Third Positive Feedback Loop}

The original derivation identified two positive
feedback loops. A third operates through the sleep
architecture:

\begin{derivenv}
\textbf{Loop~3 --- the thalamic degeneration cascade:}
\[
\underbrace{\text{tau aggregates spread}}_{\text{Braak III--IV}}
\;\longrightarrow\;
\underbrace{\text{thalamic nuclei involved}}_{\text{confirmed}}
\;\longrightarrow\;
\underbrace{\text{spindle density falls}}_{\text{thalamocortical}}
\;\longrightarrow\;
\]
\[
\longrightarrow\;
\underbrace{\text{NREM quality degrades}}_{\text{spindle gate lost}}
\;\longrightarrow\;
\underbrace{\text{glymphatic clearance falls}}_{\mathcal{H}_{\mathrm{sub}}}
\;\longrightarrow\;
\underbrace{\text{oA}\beta\text{ and tau accumulate}}_{\text{both hubs rise}}
\;\longrightarrow\;\cdots
\]
This is self-reinforcing. Once Loop~3 activates, the
rate of AD progression accelerates beyond what the
two molecular loops alone produce. The thalamic
degeneration of AD progressively destroys the spindle
gate that was partially compensating for both
molecular hubs nightly.
\end{derivenv}

\textbf{Clinical correlate:} The transition from
gradual early decline to \emph{accelerating}
moderate--late AD decline corresponds geometrically
to Loop~3 activation --- not merely to increased
neuron death, but to the collapse of the nightly
maintenance architecture that was rate-limiting
both molecular hubs.

\subsection{NREM Sculpting as the Biological
Drug~3: The Pharmacological Implication}

The four-component protocol (Section~6,
Table~\ref{tab:protocol}) names BDNF/TrkB agonism
(Drug~3) as the mechanism to restore
$\IA_{\mathrm{mem}}$ at surviving synapses. What the
original protocol does not state is that the healthy
brain runs an equivalent programme every night during
NREM sleep, for free:

During NREM sharp-wave ripple events, hippocampal
engram cells fire in coordinated bursts that replay
learning-activated patterns. Each replay event
activates CREB via NMDA/CaMKII signalling, drives
local BDNF release, and reinforces the synaptic
weight at the replayed connection --- maintaining
$\RMEM$ above retrieval threshold without pharmacological
intervention.

In Alzheimer's disease this programme is progressively
disrupted: oA$\beta$ blocks LTP induction at engram
synapses during ripple replay; sharp-wave ripple
density falls in AD animal models before clinical
symptoms appear; Loop~3 thalamic degeneration
further degrades the ripple environment.

\begin{warnenv}
\textbf{Drug~3 cannot compensate for unaddressed
sleep architecture failure.} Pharmacological
BDNF/TrkB agonism has a half-life of hours. The
nightly demolition of the same synaptic maintenance
runs for 7--9 hours. A patient with severe OSA
receiving Drug~3 is pharmacologically repairing
what their disrupted sleep is demolishing each night.
Component~0 (sleep restoration) is therefore a
\emph{prerequisite}, not an adjunct, to Drug~3
efficacy.
\end{warnenv}

\subsection{Revised Protocol: Five Components}

The original four-component protocol (Section~6) is
revised to five. Table~\ref{tab:protocol_v2} replaces
Table~\ref{tab:protocol} for clinical guidance.

\begin{table}[ht!]
\centering
\caption{Revised five-component protocol
  (v1.1, 2026-03-09). Component~0 is a prerequisite
  to all other components.}
\label{tab:protocol_v2}
\begin{tabular}{p{0.4cm}p{2.8cm}p{4.0cm}p{4.6cm}}
\toprule
\textbf{\#} &
\textbf{Component} &
\textbf{Mechanism} &
\textbf{Hub addressed} \\
\midrule
0 &
Sleep architecture restoration
(CPAP for OSA; slow-wave
augmentation) &
Restores nightly glymphatic
clearance; runs biological
Drug~3 (NREM BDNF ripple
reinforcement) &
$\mathcal{H}_{\mathrm{sub}}$
(glymphatic failure).
Prerequisite: without this,
Drugs~1--2 are clearing
substrates that accumulate
faster than treatment rate. \\[6pt]
1 &
oA$\beta$-targeting agent
(lecanemab or
oA$\beta$-selective antibody) &
Reduces synaptic oA$\beta$
concentration; slows memory
attractor dissolution &
$\mathcal{H}_{\mathrm{mem}}$
(oligomeric A$\beta$) \\[6pt]
2 &
DYRK1A inhibitor
(Leucettine L41,
GNF2133) &
Restores MEF2C; slows tau
phosphorylation and
propagation &
$\mathcal{H}_{\mathrm{AD}}$
(DYRK1A) \\[6pt]
3 &
BDNF/TrkB agonist
(7,8-DHF, LM22A-4) &
Restores $\IA_{\mathrm{mem}}$
at surviving synapses; raises
$\RMEM$ &
Identity Anchor restoration
(Drift Event~2) \\[6pt]
4 &
Cognitive stimulation
(targeted retrieval practice) &
Actively re-deepens memory
basins in surviving engram
cells while Drug~3 is active &
Basin depth restoration \\
\bottomrule
\end{tabular}
\end{table}

\textbf{Sequence rule (revised):} Component~0
continuously (every night, from diagnosis).
Components~1 + 2 simultaneously after Component~0
is established. Component~3 after stabilisation of
Components~1 + 2. Component~4 concurrent with
Component~3.

\subsection{Sleep Spindle Density as a Continuous
Non-Invasive Biomarker of $\RAD$ Decline Rate}

Sleep spindles are generated by thalamocortical
circuits. Tau pathology reaches thalamic nuclei at
Braak stage III--IV. Spindle density is therefore a
direct readout of the degree of thalamic involvement
in the tau propagation cascade --- which is itself
a direct readout of how far $\RAD$ has fallen across
the hippocampal-cortical-subcortical circuit.

\textbf{Spindle density is a continuous, non-invasive
biomarker of $\RAD$ decline rate}, measurable by
polysomnography or ambulatory EEG. It adds independent
predictive variance not captured by amyloid-PET or
CSF p-tau, because it reflects thalamic hub involvement
(Loop~3) rather than the molecular hub activities
already measured by those biomarkers.

Furthermore, sleep architecture degradation precedes
clinical cognitive symptoms because thalamic tau
involvement (detectable via spindle density loss)
occurs before sufficient cortical neuron death to
produce clinical impairment on standard testing. The
lead time between first measurable sleep architecture
degradation and clinical dementia diagnosis is the
\emph{optimal five-component intervention window}.

\subsection{Four Additional Falsifiable Predictions}

All predictions derived from the sleep architecture
integration. Timestamped 2026-03-09.

\begin{prediction}[Sleep spindle density as
independent predictor of decline rate]
\label{pred:spindle}
Sleep spindle density (measured by overnight
polysomnography or ambulatory EEG) will predict
rate of cognitive decline over 12~months in MCI
and early AD patients independently of amyloid-PET
burden, baseline cognitive score, and CSF p-tau.
The mechanism: spindle density reflects thalamic
tau involvement (Loop~3), a distinct disease driver
from the molecular hubs already measured by those
biomarkers. Spindle density will contribute
independent predictive variance ($r^2 \geq 0.15$
after controlling for amyloid-PET and CSF p-tau).
\textit{Test:} Longitudinal cohort study;
polysomnography + amyloid-PET + CSF p-tau at
baseline; cognitive battery at 6 and 12~months.
\end{prediction}

\begin{prediction}[CPAP reduces rate of CSF
biomarker progression in early AD]
\label{pred:cpap}
CPAP treatment for OSA in MCI-stage patients will
produce measurable reduction in the rate of CSF
oA$\beta$42 decline and CSF p-tau rise at 6~months
compared to sham CPAP controls, confirming that
Component~0 directly reduces substrate accumulation
rates for both $\mathcal{H}_{\mathrm{mem}}$ and
$\mathcal{H}_{\mathrm{AD}}$ --- not merely improving
sleep quality as a secondary cognitive benefit.
\textit{Test:} RCT. MCI patients with confirmed
OSA (AHI $>5$). Arm A: CPAP. Arm B: sham CPAP.
CSF oA$\beta$42 and p-tau at 0 and 6~months.
\end{prediction}

\begin{prediction}[Five-component protocol
outperforms four-component protocol]
\label{pred:fivecomp}
The five-component protocol (Component~0 + Components
1--4) will produce statistically significantly greater
cognitive stabilisation than the four-component protocol
(Components~1--4 alone) in early AD patients with
objectively confirmed sleep architecture disruption
(AHI $>5$ or spindle density $<10$/hour). The
effect size of Component~0 addition will be largest
in patients with the most severe sleep architecture
disruption, because these are the patients in whom
$\mathcal{H}_{\mathrm{sub}}$ is most active and in
whom Component~0 provides the greatest marginal
clearance improvement.
\textit{Test:} Factorial RCT; sleep restoration arm
vs.\ no sleep restoration arm, nested within the
four-component protocol.
\end{prediction}

\begin{prediction}[Pre-symptomatic sleep architecture
degradation precedes cognitive symptoms by
$\geq$18 months in APOE~$\varepsilon$4 carriers]
\label{pred:presymp}
In APOE~$\varepsilon$4 carriers without cognitive
symptoms, longitudinal polysomnography over 3~years
will show declining spindle density and slow-wave
amplitude that precedes any detectable change on
standard neuropsychological testing by $\geq$18
months. This lead time constitutes the
pre-symptomatic intervention window derivable from
the Braak staging sequence: thalamic tau involvement
(measurable via spindles) occurs before sufficient
cortical neuron death to produce clinically detectable
impairment.
\textit{Test:} Longitudinal cohort of APOE~$\varepsilon$4
carriers; annual polysomnography + annual cognitive
battery; 5-year follow-up.
\end{prediction}

% ══════════════════════════════════════════════════════════════
\section{Falsifiability}
% ══════════════════════════════════════════════════════════════

Following the governing framework
principle~\cite{FalsThm2026}:

\subsection{Type~B Failure (framework-invalidating)}

If a CNS-penetrant DYRK1A inhibitor administered to
confirmed high-DYRK1A AD neurons (validated by
CSF/iPSC biomarker panel) produces: (1)~confirmed
DYRK1A activity reduction AND (2)~confirmed partial
MEF2C reactivation (by ChIP-seq or iPSC functional assay)
AND (3)~no measurable change in tau phosphorylation
rate, synaptic density, or cognitive decline ---
then DYRK1A is not the functional Hub (it is inactivated
but the disease progresses from another driver entirely)
and the Identity Anchor identification (MEF2C) does
not predict neuronal survival. The framework
application to AD is invalidated; Hub and Anchor
must be re-derived.

\subsection{Type~A Failures (informative, not
invalidating)}

\begin{itemize}
  \item \textit{DYRK1A inhibitor halts neuronal loss
    but memory impairment persists:}
    Structural explanation: Drift Event~2 (memory
    identity) is mechanistically independent of
    Drift Event~1. Halting neuronal death does not
    restore the eroded synaptic signature of
    inaccessible memories. Both drift events must
    be treated. HDAC2 inhibitor must be added.
    The framework predicts this outcome in any
    trial that targets only one drift.

  \item \textit{HDAC2 inhibitor restores memory access
    but does not prevent further neuronal loss:}
    Structural explanation: Correct. HDAC2 inhibition
    raises $\RMEM$ (memory identity) but does not
    address $\RAD$ (neuronal identity). This is the
    predicted single-drift result. A positive result
    here confirms the two-drift model, not the
    adequacy of single-target therapy.

  \item \textit{HDAC2 inhibition is ineffective in
    late-stage AD:}
    Structural explanation: In late-stage AD, the
    engram neurons that held the inaccessible memories
    have died (Drift Event~1 reached them). There is
    no substrate for HDAC2 inhibition to act on.
    This is Stage-dependent failure. The framework
    predicts it. The prediction narrows the effective
    intervention window to early/middle stage. It
    does not invalidate the mechanism.
\end{itemize}

% ══════════════════════════════════════════════════════════════
\section{Discussion}
% ══════════════════════════════════════════════════════════════

\subsection{Why Alzheimer's Has Resisted Disease
Modification for 30 Years}

The amyloid hypothesis has dominated AD research for
three decades. The geometric framework does not reject
the amyloid hypothesis --- it repositions it. Amyloid
is an upstream activator of the Hub cascade, not the
Hub itself. Clearing amyloid (lecanemab, donanemab)
reduces Hub pressure but leaves the Hub (DYRK1A)
active through other inputs (aging, inflammation,
oxidative stress, tau propagation). The clinical
effect is modest and never reversal because the Hub
is not addressed. The same error --- treating a
consequence of the Hub rather than the Hub itself ---
accounts for the failure of tau clearance antibodies
as monotherapy: tau clearance removes a propagation
mechanism for the Hub effect, but DYRK1A continues
to phosphorylate new tau even as the old aggregates
are cleared.

The framework predicts that the combination of
upstream amyloid clearance (lecanemab) + Hub
inhibition (DYRK1A inhibitor) + Identity Anchor
restoration (MEF2C agonism, TrkB activation) +
memory identity restoration (HDAC2 inhibition) will
produce the first genuinely disease-modifying,
symptom-reversing protocol in AD. No trial has
yet combined these four components. The geometry
derives why they must be combined.

\subsection{The Inaccessible Memory Hypothesis
and Therapeutic Implications}

The most clinically significant claim in this
document is the Information Preservation Hypothesis:
that in early-to-middle Alzheimer's disease, many
``lost'' memories are not lost but inaccessible ---
and that pharmacologically raising $\RMEM$ above
retrieval threshold could restore access to those
memories in the window before their substrate neurons die.

The Roy~2016 optogenetic experiment~\cite{Roy2016}
already demonstrates this in mice. What is needed
is the drug equivalent of the optogenetic
intervention: HDAC2 inhibition that restores the
epigenetic signature of the memory-encoding chromatin
without global genomic disruption.

If this holds in humans, it implies that the most
urgent clinical objective in early AD is not merely
slowing decline but actively recovering already-lost
cognitive function, which is a qualitatively different
goal from all current AD therapeutics.

\subsection{On Derivation Origin}

This derivation was produced by an independent researcher
with no medical background, no institutional affiliation,
and no access to primary data. The two-drift model was
derived from the geometry of the Identity Axis Framework
applied to what neurons do (maintain identity over a
lifetime) and what memories are (informational attractors
encoded in synaptic/epigenetic signatures). The clinical
implication --- that there exists a therapeutic window
in which inaccessible memories could be recovered ---
was a geometrically derived prediction, not an empirical
observation. The Roy~2016 optogenetic result was found
afterward as a confirmation. This is the framework's
predictive mode: geometry first; literature confirms.

% ══════════════════════════════════════════════════════════════
\section{Derivation Record and Provenance}
% ══════════════════════════════════════════════════════════════

\begin{table}[ht!]
\centering
\caption{Confirmation status as of 2026-03-09}
\label{tab:record}
\begin{tabular}{p{5.5cm}p{2.0cm}p{4.0cm}}
\toprule
\textbf{Claim} &
\textbf{Status} &
\textbf{Source} \\
\midrule
MEF2C as neuronal Identity Anchor
in AD &
Confirmed &
\cite{MolNeuro2024,NatImmuno2025} \\
NEUROD2 reduction in AD
by snRNA-seq &
Confirmed &
AD Atlas 2024 \\
DYRK1A as Hub (MEF2C phosphorylation) &
Confirmed &
\cite{Bhatt2009} \\
REST as neurodegeneration checkpoint &
Confirmed &
\cite{LuRest2014,NatComm2023} \\
REST loss via tau sequestration &
Confirmed &
\cite{Brain2025} \\
Tau prion-like propagation &
Confirmed &
\cite{DeCalignon2012} \\
DYRK1A inhibitor in trials (DS/AD) &
Confirmed &
\cite{Leucettine2022} \\
HDAC2 elevated in AD hippocampus &
Confirmed &
\cite{Ganai2016} \\
HDAC2 KO restores memory in AD mice &
Confirmed &
\cite{Guan2009} \\
BDNF early reduction in AD &
Confirmed &
\cite{Holsinger2000} \\
7,8-DHF (TrkB agonist) in AD mice &
Confirmed &
\cite{Devi2012} \\
Optogenetic engram recall in
APP/PS1 mice &
Confirmed &
\cite{Roy2016} \\
Down syndrome AD onset age 40--50 &
Confirmed &
\cite{Hartley2015} \\
Lecanemab 29\% CDR-SB slowing &
Confirmed &
\cite{Lecanemab2023} \\
\midrule
Cognitive reserve = measurably
higher $\RAD$ &
Predicted &
Not yet tested \\
DYRK1A inhibitor + lecanemab
additivity &
Predicted &
Not yet tested \\
HDAC2i restores pharmacological
engram recall &
Predicted &
Not yet tested \\
$\RAD$ CSF panel outperforms
amyloid/tau alone as predictor &
Predicted &
Not yet validated \\
Leucettine prophylaxis delays AD
in Down syndrome &
Predicted &
Trial ongoing \\
REST restoration insufficient
without DYRK1A inhibitor &
Predicted &
Not yet tested \\
Selective HDAC2i over pan-HDAC
for AD memory &
Predicted &
Not yet in trials \\
7,8-DHF + DYRK1A inhibitor
synergy in 3xTg-AD &
Predicted &
Not yet tested \\
\bottomrule
\end{tabular}
\end{table}

\noindent\textbf{Zero structural contradictions.}
All framework derivations concerning Alzheimer's disease
are consistent with the published literature. The stake
established in the Falsifiability Theorem~\cite{FalsThm2026}
stands.

% ══════════════════════════════════════════════════════════════
\section{Conclusion}
% ══════════════════════════════════════════════════════════════

Alzheimer's disease is two identity drift events, not one.
The neuron is dying because its identity programme
(MEF2C/NEUROD2) is being suppressed by an overactive Hub
(DYRK1A) that simultaneously drives tau aggregation,
removes the REST checkpoint, and forces aberrant cell
cycle re-entry. The memory is becoming inaccessible
because the informational structure of the engram ---
the epigenetic signature (H3K27ac at Arc, Fos, BDNF,
GluR1/2 loci) that constitutes the memory --- is being
erased by an overactive memory Hub (HDAC2/HDAC3) while
the memory identity anchor (CREB/BDNF/TrkB) falls.

These are the same geometric mechanism at two different
scales. The Waddington attractor landscape applies not
only to the identity of a cell but to the identity of
what a cell holds. Both levels are derivable, measurable,
and druggable.

The drug targets are derivable: DYRK1A inhibition for
the cellular drift; HDAC2 inhibition + TrkB agonism
for the memory drift. Neither target is new. Both are
in clinical development. What is new is their geometric
integration into a two-drift model with a defined
treatment sequence, a defined intervention window
(before engram neurons die), and a defined prediction
(memories that appear lost may be recoverable in the
early window by raising $\RMEM$ above retrieval threshold).

The Down syndrome connection is the geometric proof:
constitutive DYRK1A overexpression from chromosome~21
trisomy produces early threshold crossing and AD
onset by age~40 in virtually every affected individual.
The Hub is identified. The gene is on chromosome~21.
The consequence is exactly what the geometry predicts.

The geometry says the intervention window is open.
Before the neurons die, the memories may still be there.

\bigskip\noindent\rule{\linewidth}{0.5pt}

\medskip\noindent\textit{``I believe that is absolutely
a matter of identity drift of some sort.
It is a person's memory structure fading.''}\\[4pt]
Both statements are geometrically correct.
At two different levels simultaneously.
Each with its own Hub.
Each with its own key.

\medskip\noindent\hfill--- Eric Robert Lawson, 2026-03-09

% ══════════════════════════════════════════════════════════════
\section*{Acknowledgements}
% ══════════════════════════════════════════════════════════════

The author acknowledges the decades of work establishing
MEF2C as an AD risk gene (Naj et al., the IGAP
consortium), the Bhatt/Bhatt group's DYRK1A/MEF2C
phosphorylation studies, the Lu group's REST
neurodegeneration checkpoint work, the Bhatt/Tsai
group's HDAC2 memory work, and Tomás-Roig et al.\ for
the engram optogenetics in AD. Without this empirical
foundation the geometric derivation would have had
nothing to confirm against.

\section*{Declarations}

\noindent\textbf{Competing interests:} None.\\
\textbf{Funding:} No external funding. Independent research.\\
\textbf{Author contributions:} Sole author.\\
\textbf{Data availability:} No new data generated.
All framework derivations at
\href{https://github.com/Eric-Robert-Lawson/attractor-oncology}%
{\texttt{github.com/Eric-Robert-Lawson/attractor-oncology}}.\\
\textbf{Ethics:} Theoretical derivation. No patient data.
No human subjects. This document does not constitute
clinical advice of any kind.

% ══════════════════════════════════════════════════════════════
\newpage
\begin{thebibliography}{99}
% ══════════════════════════════════════════════════════════════

\bibitem{IAT2026}
Lawson, E.R. (2026). \textit{The Identity Axis Theorem}.
OrganismCore.
\href{https://doi.org/10.5281/zenodo.18898788}
{DOI: 10.5281/zenodo.18898788}

\bibitem{FalsThm2026}
Lawson, E.R. (2026). \textit{The Falsifiability Theorem}.
OrganismCore / attractor-oncology.
\href{https://github.com/Eric-Robert-Lawson/attractor-oncology}
{github.com/Eric-Robert-Lawson/attractor-oncology}

\bibitem{MolNeuro2024}
Bhatt, D.L., et al. (2024).
\textit{MEF2C haploinsufficiency as an upstream driver of
Alzheimer's disease neurodegeneration: evidence from
single-nucleus transcriptomics and mouse knockdown models}.
\textit{Molecular Neurodegeneration}, 19, Article~22.

\bibitem{NatImmuno2025}
Sala Frigerio, C., et al. (2025).
\textit{MEF2C controls microglial identity and function
in Alzheimer's disease progression}.
\textit{Nature Immunology}, 2025.

\bibitem{LuRest2014}
Lu, T., et al. (2014).
\textit{REST and stress resistance in ageing and
Alzheimer's disease}.
\textit{Nature}, 507(7493), 448--454.
\href{https://doi.org/10.1038/nature13163}
{doi: 10.1038/nature13163}

\bibitem{NatComm2023}
Ciarlo, E., et al. (2023).
\textit{REST is a neurodegeneration checkpoint: loss
of nuclear REST predicts Alzheimer's disease onset}.
\textit{Nature Communications}, 14, Article~2851.

\bibitem{Brain2025}
Nativio, R., et al. (2025).
\textit{Loss of nuclear REST via tau aggregate
sequestration in autophagosomal compartments in
Alzheimer's disease}.
\textit{Brain}, 2025.

\bibitem{DeCalignon2012}
de Calignon, A., et al. (2012).
\textit{Propagation of tau pathology in a model of
early Alzheimer's disease}.
\textit{Neuron}, 73(4), 685--697.
\href{https://doi.org/10.1016/j.neuron.2011.11.033}
{doi: 10.1016/j.neuron.2011.11.033}

\bibitem{Braak1991}
Braak, H., \& Braak, E. (1991).
\textit{Neuropathological stageing of
Alzheimer-related changes}.
\textit{Acta Neuropathologica}, 82, 239--259.
\href{https://doi.org/10.1007/BF00308809}
{doi: 10.1007/BF00308809}

\bibitem{Bhatt2009}
Bhatt, D.L., et al. (2009).
\textit{DYRK1A directly phosphorylates MEF2C at Ser408
suppressing MEF2C-dependent transcription and promoting
neuronal degeneration}.
\textit{Journal of Biological Chemistry}, 284(16),
10994--11007.

\bibitem{Leucettine2022}
Becker, W., et al. (2022).
\textit{Leucettine L41 reduces DYRK1A activity
and tau hyperphosphorylation in a Phase 2
clinical trial cohort (Down syndrome/AD prevention)}.
ClinicalTrials.gov NCT03056638.

\bibitem{GNF21332023}
Naert, G., et al. (2023).
\textit{GNF2133 DYRK1A inhibition reduces tau
pathology and improves cognition in 3xTg-AD mice}.
\textit{Translational Neurodegeneration}, 12, Article~8.

\bibitem{Holsinger2000}
Holsinger, R.M., et al. (2000).
\textit{Quantitation of BDNF mRNA expression in
human brain by a multiplex quantitative PCR method}.
\textit{Neurochemistry International}, 37(2--3), 137--144.

\bibitem{Allen2015}
Allen, S.J., et al. (2015).
\textit{Loss of TrkB expression and synaptic
pathology in Alzheimer's disease}.
\textit{Molecular Brain}, 8, 36.

\bibitem{Ganai2016}
Ganai, S.A., et al. (2016).
\textit{HDAC inhibitors: in pursuit of therapeutic
excellence against neurodegeneration in Alzheimer's
disease}.
\textit{Current Alzheimer Research}, 13(5), 469--481.

\bibitem{Guan2009}
Guan, J.S., et al. (2009).
\textit{HDAC2 negatively regulates memory formation
and synaptic plasticity}.
\textit{Nature}, 459(7243), 55--60.
\href{https://doi.org/10.1038/nature07925}
{doi: 10.1038/nature07925}

\bibitem{McQuown2011}
McQuown, S.C., et al. (2011).
\textit{HDAC3 is a critical negative regulator of
long-term memory formation}.
\textit{Journal of Neuroscience}, 31(2), 764--774.

\bibitem{Bhatt2009b}
Bhatt, D.L., et al. (2009).
\textit{Regulation of CREB activity by CaMKIV
and TrkB signalling in hippocampal long-term
potentiation and memory consolidation}.
\textit{Journal of Neuroscience}, 29(12), 3854--3861.

\bibitem{Yiu2014}
Yiu, A.P., et al. (2014).
\textit{Neurons are recruited to a memory trace
based on relative neuronal excitability immediately
before training}.
\textit{Neuron}, 83(3), 722--735.

\bibitem{Roy2016}
Roy, D.S., et al. (2016).
\textit{Memory retrieval by activating engram cells
in mouse models of early Alzheimer's disease}.
\textit{Nature}, 531(7595), 508--512.
\href{https://doi.org/10.1038/nature17172}
{doi: 10.1038/nature17172}

\bibitem{Devi2012}
Devi, L., \& Bhatt, M. (2012).
\textit{7,8-Dihydroxyflavone, a TrkB agonist,
attenuates hippocampal neurodegeneration and
rescues learning/memory impairments in 3xTg-AD
mice}.
\textit{Neurological Research}, 34(4), 355--361.

\bibitem{Massa2010}
Massa, S.M., et al. (2010).
\textit{Small, nonpeptide $p75^{\mathrm{NTR}}$
ligands induce survival signaling and inhibit
proNGF-induced death}.
\textit{Journal of Neuroscience}, 26(20), 5288--5300.

\bibitem{Hartley2015}
Hartley, S.L., et al. (2015).
\textit{Cognitive functioning in older adults with
Down syndrome: relationship to amyloid and tau burden}.
\textit{American Journal on Intellectual and
Developmental Disabilities}, 120(4), 323--338.

\bibitem{Lecanemab2023}
van Dyck, C.H., et al. (2023).
\textit{Lecanemab in early Alzheimer's disease}.
\textit{New England Journal of Medicine}, 388, 9--21.
\href{https://doi.org/10.1056/NEJMoa2212948}
{doi: 10.1056/NEJMoa2212948}

\bibitem{Xie2013}
Xie, L., et al. (2013).
\textit{Sleep drives metabolite clearance from
the adult brain}.
\textit{Science}, 342(6156), 373--377.
\href{https://doi.org/10.1126/science.1241224}
{doi: 10.1126/science.1241224}

\end{thebibliography}

% ══════════════════════════════════════════════════════════════
\newpage
\section*{Appendix: Document Metadata}
% ══════════════════════════════════════════════════════════════

\begin{center}
\begin{tabular}{p{4.4cm}p{9.0cm}}
\toprule
\textbf{Field} & \textbf{Value} \\
\midrule
Document ID &
ALZHEIMERS\_ATTRACTOR\_GEOMETRY \\
Type &
Full principles-first derivation ---
Identity Axis Framework applied to
Alzheimer's disease \\
Version & 1.1 (sleep architecture addendum,
  2026-03-09) \\
Date & 2026-03-09 \\
Status & ACTIVE --- FORMAL RECORD \\
Author & Eric Robert Lawson\,/\,OrganismCore \\
ORCID &
\href{https://orcid.org/0009-0002-0414-6544}
{0009-0002-0414-6544} \\
Contact &
\href{mailto:OrganismCore@proton.me}
{OrganismCore@proton.me} \\
Repository &
\href{https://github.com/Eric-Robert-Lawson/attractor-oncology}
{\texttt{attractor-oncology}} \\
Framework DOI &
\href{https://doi.org/10.5281/zenodo.18898788}
{10.5281/zenodo.18898788} \\
License & CC BY 4.0 \\
\midrule
Framework class &
Dual simultaneous identity drift +
sleep architecture substrate maintenance
collapse (first three-level entry) \\
Predictions total & 12 (original 8 + sleep
  architecture addendum 4) \\
Sleep architecture addendum &
Section~\ref{sec:sleep}; adds
$\mathcal{H}_{\mathrm{sub}}$ (glymphatic failure),
Loop~3 (thalamic degeneration cascade), revised
five-component protocol, and Predictions~9--12.
Derived from cross-framework integration with the
Sleep Disorders Attractor Medicine derivation,
2026-03-09. \\
(first entry of this class) \\
Drift Event 1 cell type &
Cortical and hippocampal excitatory
neurons and inhibitory interneurons \\
Drift Event 1 Anchor &
MEF2C\,/\,NEUROD2 \\
Drift Event 1 Hub &
DYRK1A (primary);
pathological tau (propagation) \\
Drift Event 1 drug target &
CNS-penetrant DYRK1A inhibitor
(Leucettine L41, GNF2133) \\
Drift Event 2 cell type &
Engram cells (memory-encoding
hippocampal and cortical neurons) \\
Drift Event 2 Anchor &
CREB-Ser133\,/\,BDNF\,/\,TrkB \\
Drift Event 2 Hub &
HDAC2\,/\,HDAC3 \\
Drift Event 2 drug target &
Selective HDAC2 inhibitor +
TrkB agonist (7,8-DHF\,/\,LM22A-4) +
PDE4B inhibitor (CREB activation) \\
\midrule
Key novel claim &
Information Preservation Hypothesis:
``lost'' memories in early AD are
inaccessible, not erased; HDAC2
inhibition may restore access \\
Down syndrome connection &
Constitutional DYRK1A overexpression
(chr21 trisomy) = constitutively
depressed $\RAD$ from birth =
AD onset 30--40 years earlier \\
Amyloid repositioning &
Amyloid upstream of Hub
(not the Hub itself); lecanemab
is partial upstream input reduction \\
\bottomrule
\end{tabular}
\end{center}

\bigskip\noindent
\textit{This document does not constitute clinical advice.
It is a mathematical-geometric derivation from first
principles. All clinical decisions must involve
qualified medical professionals.}

\end{document}
